\documentclass[letterpaper,11pt]{article}
\usepackage[top=0.8in, bottom=0.8in, left=0.9in, right=0.9in]{geometry}
\usepackage[hidelinks]{hyperref}
\usepackage{lmodern}
\usepackage{parskip}
\usepackage{microtype}
\setlength{\parskip}{6pt}
\setlength{\parindent}{0pt}
\begin{document}
\begin{flushleft}
\textbf{\large Ajay Venkatesh} \\
Brooklyn, NY \\
+1 516 585 9013 \\
\href{mailto:av3855@nyu.edu}{av3855@nyu.edu}
\end{flushleft}
\vspace{10pt}
May 21, 2026
\vspace{6pt}

Hiring Manager \\
CalAmp

\vspace{10pt}

Dear Hiring Manager,

My name is Ajay Venkatesh. I'm finishing my M.S. in Computer Engineering at NYU, with production software experience at PwC in Bengaluru, a summer SDE internship at Amazon, and ongoing on-campus work at NYU's Novel AI Technologies. I'm writing about the Software Engineer II role on the CalAmp Telematics Cloud (CTC) team.

The posting maps cleanly onto the work I've already shipped. At Amazon I built a high-throughput, event-driven processing system on \textbf{AWS ECS Fargate, SQS, SNS} that decoupled distributed message handling and materially reduced response latency under variable load. I engineered backend services in \textbf{Java} (Coral, Smithy, Dagger, CBOR) with low-level design patterns and unit / integration testing through \textbf{Mockito}, all shipped under code-review discipline. The same posture, scalable microservices on cloud-native infrastructure processing operational data, is the shape of the CTC work the JD describes.

On performance and reliability: I provisioned \textbf{Infrastructure-as-Code} on \textbf{AWS CDK} (Lambda, API Gateway, DynamoDB, IAM, CloudWatch) with multi-stage, multi-region CI/CD and rollback controls. I'm comfortable diagnosing bottlenecks through logs, metrics, and distributed traces, and treat automated tests as part of the build rather than an afterthought. At PwC I owned a distributed Microsoft Azure backend with event-driven processing and fault-tolerant retries, where production stability was non-negotiable.

Stack-wise: \textbf{Java}, \textbf{Python}, \textbf{Kafka}, \textbf{Kubernetes}, \textbf{Docker}, \textbf{PostgreSQL}, \textbf{AWS} services across CDK and serverless, plus distributed systems debugging at scale. I'm comfortable working in a dynamic environment, contributing to architecture decisions, and translating product requirements into production-ready systems.

What pulls me toward CalAmp is the layered engineering. Connected intelligence and real-time IoT tracking sit at a layer where cloud-native software has to behave correctly against real-world devices, networks, and operational constraints, and where bugs translate into real consequences for the assets being monitored. That accountability is the environment I want to grow in.

I'd welcome the chance to discuss further. Thanks for considering my application.

\vspace{12pt}
Sincerely, \\
Ajay Venkatesh

\end{document}